\documentclass[12pt,a4paper]{article}
\usepackage[latin1]{inputenc}
\usepackage[spanish]{babel}
\usepackage{graphicx,amssymb,latexsym,amsmath,verbatim}
\usepackage{sans}
\usepackage{graphicx,amssymb,latexsym,amsmath}
\renewcommand{\thepage}{}
\renewcommand{\baselinestretch}{1}
\setlength{\parindent}{2em} \setlength{\textwidth}{18cm}
\setlength{\textheight}{24cm} \setlength{\topmargin}{-1cm}
\setlength{\oddsidemargin}{-1cm}
%\hyphenation{va-rie-da-des
%va-rie-dad pa-ra-le-lo a-si-mi-la a-pro-ve-cha-bles u-ni-for-me
%u-ni-for-me-men-te si-guien-te si-guien-tes re-su-mien-do
%ha-bi-tua-les}
%%%%%%%%%%%%%%%%%%%%%%%%%%%%%%%%%%%%%%%%%%%%%%%%%
%% Para empezar una interrogativa: ?`  %%%%%%%%%%
%%%%%%%%%%%%%%%%%%%%%%%%%%%%%%%%%%%%%%%%%%%%%%%%%
%% Para partir una matriz en cajas:         %%%%%
%% \[ \left(\begin{array}{c|c} A&B          %%%%%
%% \\ \cline{1-2} C&D\end{array}\right) \]  %%%%%
%%%%%%%%%%%%%%%%%%%%%%%%%%%%%%%%%%%%%%%%%%%%%%%%%


%SETS OF NUMBERS
%\newcommand{\R}{\mbox{${\bf R}$}}                  %reals
\newcommand{\R}{\mbox{${\mathbb R}$}}               %reals
%\newcommand{\C}{\mbox{${\bf C}$}}                  %complex
\newcommand{\C}{\mbox{${\mathbb C}$}}
%\newcommand{\N}{\mbox{${\bf N}$}}                  %natural numbers
\newcommand{\N}{\mbox{${\mathbb N}$}}
%\newcommand{\Z}{\mbox{${\bf Z}$}}                  %integers
\newcommand{\Z}{\mbox{${\mathbb Z}$}}
%\newcommand{\Q}{\mbox{${\bf Q}$}}                  %rationals
\newcommand{\Q}{\mbox{${\mathbb Q}$}}
\newcommand{\gp}{\mbox{${\bf F}_{p}$}}             % Fp
\newcommand{\pea}{\mbox{$ {\bf Z}_{p}$}}           %p-adic integers
\newcommand{\zn}{\mbox{${\bf Z}/{n}$}}             %Z/n
\newcommand{\zp}{\mbox{${\bf Z}/{p}$}}             %Z/p
\newcommand{\Zh}{\mbox{$\hat{{\bf Z}}$}}           %Z^
\newcommand{\xp}{\mbox{${\bf F}_{p}^{*}$}}         %units Fp

%GREEK LETTERS
\newcommand{{\ff}}{\mbox{$\varphi$}}                   %phi
\newcommand{\fp}{\mbox{$\varphi_{p}$}}               %phi_p
\newcommand{\al}{\mbox{$\alpha$}}                    %alpha
\newcommand{\be}{\mbox{$\beta$}}                     %beta
\newcommand{\px}{\mbox{$\varepsilon$}}               %epsilon
\newcommand{\Ga}{\mbox{$\Gamma$}}                    %capital gamma
\newcommand{\ga}{\mbox{$\gamma$}}                    %small gamma

%LOGIC SYMBOLS
\newcommand{\ex}{\mbox{$\exists$}}                   %there is
\newcommand{\fa}{\mbox{$\forall$}}                   %forall
\newcommand{\stf}{\mbox{$\models$}}                  %satisfies
\newcommand{\y}{\mbox{$\wedge$}}                     %and
%\newcommand{\o}{\mbox{$\vee$}}                       %or
\newcommand{\si}{\mbox{$\Longrightarrow$}}           %implies
\newcommand{\ekiv}{\mbox{$\Longleftrightarrow$}}     %equiv

%OTHERS
\newcommand{\dm}{\mbox{$\preceq$}}                   %dominates




\begin{document}

\small

\noindent{\sf Universidad Aut\'onoma de Madrid  \hfill \'Algebra.
Curso 2013--14} \vspace{0.1cm} \hrule \vspace{0.1cm} \noindent{\sf
Ingenier\'ia Inform\'atica} \vspace{0.1cm}
\begin{center}
{\bfseries  \'ALGEBRA. HOJA 6.}
\end{center}

\begin{enumerate}
%\renewcommand{\theenumi}{\bf\arabic{enumi})\ }
\renewcommand{\labelenumi}{\bf\arabic{enumi})\ }
%\renewcommand{\theenumii}{\alph{enumii})\ }
\renewcommand{\labelenumii}{\alph{enumii})\ }
\medskip



\item Calcula los siguientes determinantes:
\begin{align*}
 & a) \
 \begin{vmatrix}
 4 & -2 & 5 & 1 \\
 -4 & 1 & 0 & -1 \\
 2 & 1 & -1 & 1 \\
 1 & 0 & 1 & -2
 \end{vmatrix}
 && b) \
  \begin{vmatrix}
 -4 & 1 & 1 & 1 & 1 \\
 1 & -4 & 1 & 1 & 1 \\
 1 & 1 & -4 & 1 & 1 \\
 1 & 1 & 1 & -4 & 1 \\
 1 & 1 & 1 & 1 & -4
 \end{vmatrix}
 && c) \
 \begin{vmatrix}
 x & a & a & a \\
 a & x & a & a  \\
 a & a & x & a \\
 a & a & a & x
 \end{vmatrix}
 && d) \
 \begin{vmatrix}
 1 & 2 & 3 & 4 & 5 \\
 2 & 2 & 3 & 4 & 5 \\
 3 & 3 & 3 & 4 & 5 \\
 4 & 4 & 4 & 4 & 5 \\
 5 & 5 & 5 & 5 & 5
 \end{vmatrix}  \\
 & e) \
 \begin{vmatrix}
 a_{11} & a_{12} & a_{13} & \cdots & a_{1n} \\
 0 & a_{22} & a_{23} & \cdots & a_{2n} \\
 0 & 0 & a_{33} & \cdots & a_{3n} \\
 \vdots & \vdots & \vdots & \ddots & \vdots \\
 0 & 0 & 0 & \cdots & a_{nn}
 \end{vmatrix} .
\end{align*}

\noindent
\medskip\textbf{Soluci\'on: } a) $63$;
b) $0$;
c) $(x-a)^3 (3a + x)$;
d) $5$;
e) $a_{11} \cdot a_{22} \cdot a_{33} \cdots a_{nn}$.

\medskip


\item Utilizando determinantes, calcula, cuando exista, la inversa de las siguientes matrices
\begin{equation*}
 a) \
 \begin{pmatrix}
 1 & 2 & 3 & 4 \\
 2 & 3 & 5 & -2 \\
 3 & 5 & 1 & 1 \\
 1 & 1 & -5 & -7
 \end{pmatrix} \quad
 b) \
  \begin{pmatrix}
 3 & 3 & 2 \\
 4 & 3 & 3 \\
 1 & 1 & 1
 \end{pmatrix} \quad
 c) \
 \begin{pmatrix}
 a & 1 & 1  \\
 2 & -1 & 1 \\
 b & -2 & -1
 \end{pmatrix} .
\end{equation*}


\noindent
\medskip\textbf{Soluci\'on: } a) No existe;
b) $\begin{pmatrix}
      0  &   1  &  -3\\
      1  &  -1  &   1\\
      -1  &  0  &   3
\end{pmatrix}$;\\
c) $\displaystyle \frac{1}{3a + 2b - 2}
\begin{pmatrix}
 3 & -1 &   2 \\
 b + 2 &  -a - b & -a + 2  \\
 b - 4 & 2a + b & -a - 2
\end{pmatrix}$
si $3a + 2b - 2 \neq 0$; no existe si $3a + 2b - 2 =0$.


\medskip

\item Decide, usando determinantes, si la matriz de coeficientes de los siguientes sistemas de ecuaciones es invertible. Cuando lo sea, calcula la soluci\'on del sistema usando la regla de Cramer.
{\setlength\arraycolsep{.1em}
\begin{equation*}
 a) \ \left.
\begin{array}{rrrrrrrr}
x_1 &+& 2x_2 &-& 3x_3 &=& 7 \\
2 x_1  &+& x_2 &+ & x_3 &=& 6\\
a x_1  &-& x_2 &+& 3 x_3 &=& -1
\end{array}
\right\} \qquad
b) \ \left.
\begin{array}{rrrrrrrr}
2x_1 &-& 4x_2 & & &=& b \\
x_1 &-& 3x_2  & & &=& -4 \\
x_1 &+& a x_2 &-& x_3  &=& 4
\end{array}
\right\}
\end{equation*}
}

\noindent
\medskip\textbf{Soluci\'on: } a) No es invertible si $a = \frac{2}{5}$; cuando $a \neq \frac{2}{5}$ la soluci\'on es $\left( \frac{5}{5a - 2} , \frac{25a - 17}{5a - 2} , \frac{5 a - 5}{5a - 2} \right)$; \\
b) $\left( \frac{3b}{2} + 8 , \frac{b}{2} + 4 , \frac{4a}{3} + \frac{b}{2} + \frac{ab}{6} + \frac{4}{3} \right)$.
\medskip

\item Comprueba que
\[
  \begin{vmatrix}
 a_{11} & a_{12} & a_{13} & a_{14} \\
 a_{21} & a_{22} & a_{23} & a_{24} \\
 0 & 0 & a_{33} & a_{34} \\
 0 & 0 & a_{43} & a_{44}
 \end{vmatrix} =
\begin{vmatrix}
 a_{11} & a_{12} \\
 a_{21} & a_{22}
 \end{vmatrix} \cdot \begin{vmatrix}
  a_{33} & a_{34} \\
 a_{43} & a_{44}
 \end{vmatrix} .
\]


\end{enumerate}




\end{document}
